\section{Reti Logiche}
Le reti logiche sono dei circuiti che elaborano segnali logici booleani tramite la composizione dei gate logici fondamentali. Possono essere usate per calcolare funzioni o per memorizzare segnali.

\subsection{Reti logiche combinatorie}
\label{subsec:reti-comb}

Le reti logiche combinatorie modellano relazioni ingresso-uscita senza memoria: l'output dipende \emph{solo} dagli input presenti. La progettazione segue un flusso standard:

\begin{enumerate}
    \item Descrizione della funzione booleana
    \item Costruzione della tabella di verità
    \item Derivazione dell'espressione logica (forma SOP -- \textit{sum of products})
    \item Sintesi della rete di porte
\end{enumerate}

Con ``somma di prodotti'' (SOP) si intende la scrittura della funzione come OR (somma) di termini AND (prodotti), derivati direttamente dalle righe della tabella di verità in cui l'uscita vale~1.

\subsubsection*{Esempio: progettazione di un Full Adder a 1 bit}
\label{subsec:fa1-progetto}

Il \textit{Full Adder} è un circuito combinatorio che somma due bit $x$ e $y$ insieme a un riporto in ingresso $c_{\text{in}}$, producendo una somma $z$ e un riporto in uscita $c_{\text{out}}$.

\paragraph{Specifica del problema.}
\begin{itemize}
    \item \textbf{Input:} bit $x$, bit $y$, riporto in ingresso $c_{\text{in}}$
    \item \textbf{Output:} bit somma $z$, riporto in uscita $c_{\text{out}}$
\end{itemize}

\paragraph{Descrizione della funzione.}
L'obiettivo è realizzare un sommatore completo a 1 bit: dati due bit da sommare e un eventuale riporto proveniente da una somma precedente (per composizionalità in somme multi-bit), il circuito deve calcolare il bit di somma e il nuovo riporto da propagare.

\paragraph{Costruzione della tabella di verità.}
\[
\begin{array}{ccc|cc}
\toprule
x & y & c_{\text{in}} & z & c_{\text{out}} \\
\midrule
0 & 0 & 0 & 0 & 0 \\
0 & 0 & 1 & 1 & 0 \\
0 & 1 & 0 & 1 & 0 \\
0 & 1 & 1 & 0 & 1 \\
1 & 0 & 0 & 1 & 0 \\
1 & 0 & 1 & 0 & 1 \\
1 & 1 & 0 & 0 & 1 \\
1 & 1 & 1 & 1 & 1 \\
\bottomrule
\end{array}
\]

\paragraph{Derivazione della forma canonica SOP.}
Per ottenere l'espressione booleana, individuiamo nella tabella di verità tutte le righe in cui l'uscita vale~1, separatamente per~$z$ e per~$c_{\text{out}}$.

Per la somma~$z$, l'uscita vale~1 nella seconda, terza, quinta e ottava riga. Da ciascuna riga costruiamo un \emph{mintermine} (prodotto logico) delle variabili di ingresso: se un ingresso vale~0 usiamo la forma negata, se vale~1 la forma diretta. La somma (OR) di tutti i mintermini produce l'espressione per~$z$:

\begin{equation}
\label{eq:z-sop}
z = \overline{x} \cdot \overline{y} \cdot c_{\text{in}} + \overline{x} \cdot y \cdot \overline{c_{\text{in}}} + x \cdot \overline{y} \cdot \overline{c_{\text{in}}} + x \cdot y \cdot c_{\text{in}}
\end{equation}

L'Equazione~\ref{eq:z-sop} elenca tutte e sole le combinazioni di input per cui $z = 1$.

Procediamo analogamente per il riporto in uscita $c_{\text{out}}$, sommando i mintermini relativi alle righe in cui $c_{\text{out}} = 1$:

\begin{equation}
\label{eq:cout-sop}
c_{\text{out}} = x \cdot y \cdot \overline{c_{\text{in}}} + x \cdot \overline{y} \cdot c_{\text{in}} + \overline{x} \cdot y \cdot c_{\text{in}} + x \cdot y \cdot c_{\text{in}}
\end{equation}

\paragraph{Sintesi della rete di porte.}
A partire dalle espressioni SOP (Equazioni~\ref{eq:z-sop} e~\ref{eq:cout-sop}), la sintesi circuitale procede come segue:
\begin{enumerate}
    \item Creare una linea verticale per ogni segnale di ingresso ($x$, $y$, $c_{\text{in}}$) e le rispettive negazioni.
    \item Realizzare ciascun prodotto con una porta AND a tre ingressi, collegando le linee corrette (dirette o negate).
    \item Collegare le uscite delle porte AND agli ingressi di una porta OR per ciascun output.
    \item L'uscita della porta OR corrisponde al segnale di output finale ($z$ o $c_{\text{out}}$).
\end{enumerate}

\begin{figure}[htbp]
    \centering
    \includegraphics[scale=0.4]{reti_comb_fa1.png}
    \caption{Implementazione circuitale del Full Adder a 1 bit: le linee verticali rappresentano i segnali di ingresso $x$, $y$ e $c_{\text{in}}$ (con le relative negazioni); le porte AND realizzano i prodotti della forma SOP (Equazione~\ref{eq:z-sop} per la somma, Equazione~\ref{eq:cout-sop} per il riporto); le porte OR combinano i prodotti per produrre le uscite finali $z$ e $c_{\text{out}}$.}
    \label{fig:fa1}
\end{figure}

\subsection{Ottimizzazione nella logica combinatoria}

Spesso e volentieri la formula che scriviamo non è la più ottimale. Spesso e volentieri  possiamo usare dei piccoli trucchi di algebra booleana per renderla più facile.

\subsection{Verilog}
\subsection{Reti logiche sequenziali}
\subsubsection{Automi di Mealy}
\subsubsection{Automi di Moore}
\subsection{Componenti standard}
\subsubsection{Full Adder}
\subsubsection{Mux e Demux}
\subsubsection{Comparatore}
\subsubsection{Shifter}
\subsubsection{Encoder e Decoder}
